\chapter{Relaxation and dephasing}
\label{ch:relaxation-dephasing}

The Lindblad equation of \cref{ch:open-systems} is general; this
chapter specialises it to the single most important open system in
the book, a qubit coupled to a noisy environment. Two qualitatively
different things can happen to a qubit in contact with a bath. It
can \emph{lose energy} --- decay from $\ket1$ to $\ket0$ by emitting
a quantum into the environment --- which is \emph{relaxation},
characterised by the time $T_1$. Or it can \emph{lose phase
coherence} --- the relative phase between $\ket0$ and $\ket1$ can be
randomised by fluctuations of its transition frequency, without any
energy exchange at all --- which is \emph{pure dephasing},
characterised by the time $T_\phi$. The experimentally measured
coherence time $T_2$ combines the two. We derive each from the
appropriate jump operator, assemble the optical Bloch equations,
establish the relation $1/T_2 = 1/(2T_1)+1/T_\phi$, distinguish
$T_2$ from the inhomogeneous free-induction time $T_2^*$, and close
by connecting all of these rates to the environmental noise
spectral density --- the viewpoint that explains why $T_1$ probes
the bath at the qubit frequency while dephasing probes it near zero
frequency.

% --------------------------------------------------------------
\section{Energy relaxation: \texorpdfstring{$T_1$}{T1}}
\label{sec:relaxation}

\begin{phenomenon}
A qubit prepared in its excited state $\ket1$ does not stay there.
Coupled to a cold electromagnetic environment, it emits a quantum
and decays to $\ket0$; coupled to a bath at finite temperature, it
also occasionally absorbs a quantum and is promoted back to $\ket1$.
The net exponential approach of the excited-state population to its
thermal value defines the energy-relaxation time $T_1$. Relaxation
is an \emph{energy-exchanging} process: it requires the bath to be
able to accept (or supply) a quantum at the qubit frequency
$\omega_q$.
\end{phenomenon}

\begin{statement}[\textbf{Relaxation from amplitude damping}]
\label{stmt:t1}
A qubit with $\op H=-\tfrac{\hbar\omega_q}{2}\hat\sigma_z$ coupled to
a thermal bath is described by the Lindblad equation
(\cref{stmt:lindblad}) with two jump operators: emission
$\op L_\downarrow=\hat\sigma_-$ at rate $\gamma_\downarrow$ and
absorption $\op L_\uparrow=\hat\sigma_+$ at rate $\gamma_\uparrow$.
The excited-state population obeys
\begin{equation}
\label{eq:t1-population}
\dot\rho_{11} \;=\; -\gamma_\downarrow\,\rho_{11} + \gamma_\uparrow\,\rho_{00},
\end{equation}
relaxing exponentially to the steady value
$\rho_{11}^{\infty}=\gamma_\uparrow/(\gamma_\uparrow+\gamma_\downarrow)$
with the energy-relaxation rate
\begin{equation}
\label{eq:t1-rate}
\frac{1}{T_1} \;=\; \gamma_\downarrow + \gamma_\uparrow.
\end{equation}
Detailed balance fixes the ratio
$\gamma_\uparrow/\gamma_\downarrow=e^{-\hbar\omega_q/k_B T}$, so at
$T\ll\hbar\omega_q/k_B$ absorption is exponentially suppressed,
$1/T_1\approx\gamma_\downarrow$ and $\rho_{11}^\infty\approx0$.
\end{statement}

\paragraph{Derivation.}
Starting from \cref{eq:amplitude-damping-me} generalised to two
jump operators, the dissipator
$\gamma_\downarrow\mathcal D[\hat\sigma_-]+\gamma_\uparrow\mathcal D[\hat\sigma_+]$
acts on the populations as in \cref{eq:t1-population} (each jump
moves population one way, each anticommutator term removes it),
and the unitary part $-\tfrac{i}{\hbar}\comm{\op H}{\op\rho}$ leaves
populations untouched. The off-diagonal coherence picks up
$\dot\rho_{01}=(i\omega_q-\tfrac12(\gamma_\downarrow+\gamma_\uparrow))\rho_{01}$,
so even pure relaxation already damps the coherence at half the
energy-relaxation rate --- the unavoidable contribution of $T_1$ to
the coherence budget, made precise in \cref{sec:bloch-equations} below.

\paragraph{Connection to Fermi's golden rule.}
The emission rate $\gamma_\downarrow$ is exactly the spontaneous
quantum-decay rate one would compute from Fermi's golden rule
(\cref{thm:fermi}) for the qubit--bath coupling evaluated at the
transition frequency. The Markovian master equation reproduces this
rate and additionally tracks the coherence, the absorption channel,
and the approach to thermal equilibrium --- information that the
single golden-rule number does not carry.

\paragraph{Numerical scale.}
For a transmon at $\omega_q/2\pi=5\,$GHz the thermal occupancy at
$T=20\,$mK is $\bar n=(e^{\hbar\omega_q/k_BT}-1)^{-1}\sim10^{-5}$, so
the steady-state excited population is at the $10^{-5}$ level and
relaxation is essentially one-way decay. Present-day transmons
reach $T_1\sim 100\,\mu$s--$1\,$ms, set by dielectric loss, the
Purcell channel into the readout resonator, and quasiparticle
tunnelling (\cref{stmt:qp-relaxation}).

% --------------------------------------------------------------
\section{Pure dephasing: \texorpdfstring{$T_\phi$}{Tphi}}
\label{sec:dephasing}

\begin{phenomenon}
A qubit can lose phase coherence without losing any energy. If the
transition frequency $\omega_q$ fluctuates --- because of charge,
flux, or photon-number noise in the environment --- then the
relative phase accumulated between $\ket0$ and $\ket1$ becomes
random from run to run, and the ensemble-averaged coherence decays.
Because no quantum is exchanged with the bath, this \emph{pure
dephasing} can occur even when $T_1\to\infty$, and it probes the
environment at low frequency rather than at $\omega_q$.
\end{phenomenon}

\begin{statement}[\textbf{Pure dephasing from a $\hat\sigma_z$ channel}]
\label{stmt:tphi}
A Markovian frequency-noise environment is described by the
Lindblad equation with the single Hermitian jump operator
$\op L=\hat\sigma_z$ at rate $\gamma_z$,
\begin{equation}
\label{eq:dephasing-me}
\frac{\dd\op\rho}{\dd t}
= -\frac{i}{\hbar}\comm{\op H}{\op\rho}
+ \gamma_z\,\mathcal D[\hat\sigma_z]\,\op\rho.
\end{equation}
Because $\hat\sigma_z$ is diagonal, the populations are untouched
($\dot\rho_{00}=\dot\rho_{11}=0$) while the coherence decays purely
exponentially,
\begin{equation}
\label{eq:tphi-decay}
\dot\rho_{01} \;=\; \bigl(i\omega_q - 2\gamma_z\bigr)\rho_{01},
\qquad \frac{1}{T_\phi} \;=\; 2\gamma_z.
\end{equation}
The pure-dephasing rate $1/T_\phi$ is the rate at which the
off-diagonal element decays in excess of the relaxation
contribution.
\end{statement}

\paragraph{Derivation.}
For $\op L=\hat\sigma_z$ the dissipator is
$\mathcal D[\hat\sigma_z]\op\rho=\hat\sigma_z\op\rho\hat\sigma_z-\op\rho$
(using $\hat\sigma_z^2=\id$). In the
$\{\ket0,\ket1\}$ basis $\hat\sigma_z\op\rho\hat\sigma_z$ flips the
sign of the off-diagonal elements and leaves the diagonal alone, so
$\mathcal D[\hat\sigma_z]\op\rho$ has zero diagonal and
$-2\rho_{01}$ off-diagonal. Hence the populations are constant and
$\rho_{01}$ decays at $2\gamma_z$, giving \cref{eq:tphi-decay}. This
is the Markovian, exponential version of the dephasing kernel
derived microscopically by tracing out an environment in
\cref{ex:decoherence-by-partial-trace}: there the off-diagonal
element was multiplied by an overlap factor that decayed in time;
here that decay is made exponential by the Markov approximation.

\paragraph{Quasi-static versus Markovian dephasing.}
\Cref{stmt:tphi} assumes white (Markovian) frequency noise, which
gives exponential coherence decay. Real solid-state environments
are dominated by $1/f$ noise --- slow, strongly correlated
fluctuations of charge and flux --- whose low-frequency weight is
large. Slow noise is better described as \emph{quasi-static}: the
qubit frequency is essentially constant within one experimental
shot but varies from shot to shot, producing a \emph{Gaussian}
(rather than exponential) decay of the averaged coherence and a
free-induction time $T_2^*$ that can be much shorter than the
Markovian estimate. The distinction is the subject of
\cref{sec:bloch,sec:noise-psd}.

% --------------------------------------------------------------
\section{The Bloch equations and the coherence-time relations}
\label{sec:bloch-equations}

\begin{statement}[\textbf{Optical Bloch equations}]
\label{stmt:bloch-equations}
Combining relaxation (\cref{stmt:t1}) and pure dephasing
(\cref{stmt:tphi}), the components of the Bloch vector
$\vec r=(\langle\hat\sigma_x\rangle,\langle\hat\sigma_y\rangle,\langle\hat\sigma_z\rangle)$
(\cref{def:density-tls}) obey
\begin{equation}
\label{eq:bloch-equations}
\dot r_x = -\omega_q r_y - \frac{r_x}{T_2},\quad
\dot r_y = +\omega_q r_x - \frac{r_y}{T_2},\quad
\dot r_z = -\frac{r_z - r_z^\infty}{T_1},
\end{equation}
with $r_z^\infty=\rho_{11}^\infty-\rho_{00}^\infty$ the thermal
steady state. The transverse components precess at $\omega_q$ and
decay at the \emph{total} dephasing rate $1/T_2$, while the
longitudinal component relaxes at $1/T_1$. The three timescales are
related by
\begin{equation}
\label{eq:t2-relation}
\boxed{\;\frac{1}{T_2} \;=\; \frac{1}{2T_1} \;+\; \frac{1}{T_\phi}\;}
\end{equation}
\end{statement}

\paragraph{Reading the relation.}
\Cref{eq:t2-relation} says the coherence decays for two independent
reasons: relaxation, which destroys the coherence at half the
energy-decay rate (the $1/2T_1$ floor derived in
\cref{sec:relaxation}), and pure dephasing, which adds $1/T_\phi$ on
top. In the absence of pure dephasing ($T_\phi\to\infty$) the
coherence is \emph{relaxation-limited}, $T_2=2T_1$, the longest a
coherence can survive given a fixed $T_1$. Any frequency noise
shortens it.

\begin{statement}[\textbf{$T_2$ versus $T_2^*$}]
\label{stmt:t2star}
The time $T_2$ in \cref{eq:t2-relation} is the \emph{homogeneous}
(intrinsic) dephasing time, the coherence decay seen after
echo sequences remove slow noise. The \emph{inhomogeneous} or
free-induction time $T_2^*$ is what a simple Ramsey experiment
measures: it additionally includes quasi-static, shot-to-shot
frequency fluctuations (slow charge/flux drift, ensemble spread),
and satisfies
\begin{equation}
\label{eq:t2star-ineq}
T_2^* \;\leq\; T_2.
\end{equation}
$T_2^*$ is not a pure-dephasing time and should not be identified
with $T_\phi$; it is the convolution of the homogeneous decay with
the quasi-static noise distribution, often Gaussian in form.
\end{statement}

\begin{remark}[Terminology]
\label{rem:t2-terminology}
The pure-dephasing time is denoted $T_\phi$ (rate
$\Gamma_\phi=1/T_\phi$), \emph{not} $T_2^*$. The free-induction time
$T_2^*$ is the inhomogeneous dephasing time including quasi-static
detuning noise and shot-to-shot frequency variations. The three
homogeneous times are tied by \cref{eq:t2-relation}; the
inhomogeneous $T_2^*$ sits below $T_2$ by \cref{eq:t2star-ineq}.
Echo and CPMG pulse sequences (\cref{sec:experimental-basics}, in
Part~IV) refocus the slow noise and recover $T_2$ from $T_2^*$.
\end{remark}

\paragraph{Transition.}
The Bloch picture tells us \emph{how} the coherences and
populations decay, but the rates $1/T_1$ and $1/T_\phi$ are still
phenomenological inputs. The final section ties them to a single
microscopic object --- the environmental noise spectral density ---
and shows that relaxation and dephasing simply sample that spectrum
at different frequencies.

% --------------------------------------------------------------
\section{Noise spectral density and the sampling picture}
\label{sec:noise-psd}

\begin{phenomenon}
Both relaxation and dephasing arise from a fluctuating coupling
between the qubit and its environment. The strength of those
fluctuations as a function of frequency is the \emph{noise spectral
density} $S(\omega)$, the Fourier transform of the bath
autocorrelation function. Relaxation requires the environment to
absorb a quantum at the qubit frequency, so $1/T_1$ samples
$S(\omega)$ at $\omega=\omega_q$; pure dephasing arises from
slow frequency wander, so $1/T_\phi$ samples $S(\omega)$ near
$\omega=0$. The same bath, read at two different frequencies,
controls the two decay channels.
\end{phenomenon}

\begin{statement}[\textbf{Rates from the noise spectrum}~\cite{clerk2010}]
\label{stmt:noise-rates}
Write the qubit--bath coupling as
$\op H_{SB}=\hbar(\hat\sigma_z\,\hat\lambda_\parallel/2 + \hat\sigma_x\,\hat\lambda_\perp)$,
with $\hat\lambda_\parallel$ a longitudinal (frequency) noise
operator and $\hat\lambda_\perp$ a transverse (coupling) noise
operator, and let $S_\parallel(\omega)$, $S_\perp(\omega)$ be their
spectral densities. Then to leading order in the coupling
(Fermi's golden rule, \cref{thm:fermi}),
\begin{equation}
\label{eq:t1-from-noise}
\frac{1}{T_1} \;=\; S_\perp(\omega_q) + S_\perp(-\omega_q),
\end{equation}
the sum of the emission ($+\omega_q$) and absorption ($-\omega_q$)
weights of the transverse spectrum, while the pure-dephasing rate
for white longitudinal noise is
\begin{equation}
\label{eq:tphi-from-noise}
\frac{1}{T_\phi} \;=\; \tfrac12\,S_\parallel(0),
\end{equation}
the zero-frequency weight of the longitudinal spectrum. For
non-white (e.g.\ $1/f$) longitudinal noise the simple
white-noise rate \cref{eq:tphi-from-noise} is replaced by a
filter-function overlap, producing the non-exponential decay noted
in \cref{stmt:t2star}.
\end{statement}

\paragraph{Why the frequencies differ.}
The transverse coupling $\hat\sigma_x$ connects $\ket0$ and $\ket1$,
so it can drive transitions only by exchanging the energy
$\hbar\omega_q$ --- it is resonant, and samples $S_\perp$ at
$\pm\omega_q$. The longitudinal coupling $\hat\sigma_z$ commutes
with $\op H_S$: it does not cause transitions but modulates the
energy splitting, so its effect accumulates from the slow, near-DC
part of the spectrum. This frequency separation is the basis of
\emph{noise spectroscopy}: by measuring $T_1$ and $T_\phi$ (and
their dependence on dynamical-decoupling sequences) one reconstructs
$S(\omega)$ over a wide band, a routine cQED characterisation tool.

\paragraph{The fluctuation--dissipation connection.}
The thermal-equilibrium spectral density obeys the
fluctuation--dissipation relation
$S(\omega)=\hbar\omega\,[\,\coth(\hbar\omega/2k_BT)+1\,]\,\mathrm{Re}\,Z(\omega)/\ldots$,
which ties the dissipative response of the environment (an
impedance $Z(\omega)$, for an electrical bath) to the noise it
injects. Operationally, this is why a low-loss, low-temperature
environment is a quiet one: suppressing dissipation at $\omega_q$
suppresses the relaxation noise, and suppressing low-frequency loss
suppresses the dephasing noise.

% --------------------------------------------------------------
This chapter turned the abstract Lindblad equation into the
working vocabulary of qubit coherence. Relaxation comes from the
amplitude-damping channel and defines $T_1$; pure dephasing comes
from the $\hat\sigma_z$ channel and defines $T_\phi$; the optical
Bloch equations assemble them, with the total coherence time set by
$1/T_2=1/(2T_1)+1/T_\phi$ and the free-induction time $T_2^*\leq T_2$
further shortened by quasi-static noise. All the rates are
ultimately set by the environmental noise spectral density,
sampled at $\omega_q$ for relaxation and near zero for dephasing.

These are the numbers every cQED experiment reports. The next
chapter, \cref{ch:input-output}, returns to the microscopic bath
and treats it explicitly as a continuum of propagating modes,
deriving the input--output relations that connect the intracavity
field to the signals travelling on the lines --- the formalism
behind dispersive readout and the photon-counter efficiency of
Part~VI.

\clearpage
\begin{chaptersummary}

\paragraph{Relaxation $T_1$.}
The amplitude-damping channel $\hat\sigma_-$ (rate
$\gamma_\downarrow$) plus thermal absorption $\hat\sigma_+$ (rate
$\gamma_\uparrow$) gives $1/T_1=\gamma_\downarrow+\gamma_\uparrow$,
with steady-state population
$\rho_{11}^\infty=\gamma_\uparrow/(\gamma_\uparrow+\gamma_\downarrow)$
and detailed balance
$\gamma_\uparrow/\gamma_\downarrow=e^{-\hbar\omega_q/k_BT}$.

\paragraph{Pure dephasing $T_\phi$.}
The $\hat\sigma_z$ channel (rate $\gamma_z$) leaves populations
fixed and decays the coherence at $1/T_\phi=2\gamma_z$, with no
energy exchange.

\paragraph{Bloch equations and $T_2$.}
Transverse components precess at $\omega_q$ and decay at $1/T_2$;
the longitudinal component relaxes at $1/T_1$. The homogeneous
times satisfy
\begin{equation}
\frac{1}{T_2} = \frac{1}{2T_1} + \frac{1}{T_\phi},
\end{equation}
so $T_2\leq 2T_1$, with equality only when pure dephasing is
absent. The inhomogeneous free-induction time obeys $T_2^*\leq T_2$
and includes quasi-static, shot-to-shot frequency noise; it is not
the pure-dephasing time.

\paragraph{Noise spectral density.}
$1/T_1=S_\perp(\omega_q)+S_\perp(-\omega_q)$ samples transverse
noise at the qubit frequency; $1/T_\phi=\tfrac12 S_\parallel(0)$
samples longitudinal noise near DC. Measuring the two (and their
dynamical-decoupling dependence) reconstructs $S(\omega)$.

\end{chaptersummary}

% --------------------------------------------------------------
\section*{Exercises}
\addcontentsline{toc}{section}{Exercises}

\begin{exercise}[Solving the population equation]
Solve \cref{eq:t1-population} for $\rho_{11}(t)$ given
$\rho_{11}(0)=1$, and show it relaxes exponentially to
$\rho_{11}^\infty$ with time constant $T_1$.
\begin{solution}
With $\rho_{00}=1-\rho_{11}$, \cref{eq:t1-population} reads
$\dot\rho_{11}=-(\gamma_\downarrow+\gamma_\uparrow)\rho_{11}+\gamma_\uparrow$.
The solution is
$\rho_{11}(t)=\rho_{11}^\infty+(1-\rho_{11}^\infty)e^{-t/T_1}$ with
$\rho_{11}^\infty=\gamma_\uparrow/(\gamma_\uparrow+\gamma_\downarrow)$
and $1/T_1=\gamma_\downarrow+\gamma_\uparrow$. It decays
exponentially from $1$ to $\rho_{11}^\infty$.
\end{solution}
\end{exercise}

\begin{exercise}[Thermal occupancy of a transmon]
Compute the equilibrium excited-state population
$\rho_{11}^\infty$ for $\omega_q/2\pi=5\,$GHz at $T=20\,$mK and at
$T=100\,$mK. Comment on the implication for state-preparation
fidelity.
\begin{solution}
$\hbar\omega_q/k_B = (h\cdot5\times10^9)/k_B\approx 0.24\,$K. At
$20\,$mK, $\hbar\omega_q/k_BT\approx12$, so
$\rho_{11}^\infty\approx e^{-12}\approx6\times10^{-6}$. At
$100\,$mK, $\approx e^{-2.4}\approx0.09$ (using
$\rho_{11}^\infty=1/(1+e^{\hbar\omega_q/k_BT})$ gives $\approx0.08$).
Residual thermal population caps passive ground-state preparation
fidelity; at $100\,$mK an $8\%$ excited population is unacceptable,
motivating active reset.
\end{solution}
\end{exercise}

\begin{exercise}[The $\hat\sigma_z$ dissipator]
Show that $\mathcal D[\hat\sigma_z]\op\rho=\hat\sigma_z\op\rho\hat\sigma_z-\op\rho$
and evaluate its action on a general single-qubit $\op\rho$ in the
$\{\ket0,\ket1\}$ basis. Confirm \cref{eq:tphi-decay}.
\begin{solution}
Since $\hat\sigma_z^\dagger\hat\sigma_z=\id$,
$\mathcal D[\hat\sigma_z]\op\rho=\hat\sigma_z\op\rho\hat\sigma_z-\tfrac12\{\id,\op\rho\}
=\hat\sigma_z\op\rho\hat\sigma_z-\op\rho$. With
$\hat\sigma_z=\mathrm{diag}(1,-1)$ (taking $\ket0=(1,0)$),
$\hat\sigma_z\op\rho\hat\sigma_z$ has unchanged diagonal and negated
off-diagonal, so $\mathcal D[\hat\sigma_z]\op\rho$ has zero diagonal
and $-2\rho_{01}$ at the off-diagonal. Adding the unitary precession
gives $\dot\rho_{01}=(i\omega_q-2\gamma_z)\rho_{01}$, i.e.\
$1/T_\phi=2\gamma_z$.
\end{solution}
\end{exercise}

\begin{exercise}[The relaxation-limited coherence time]
Show that pure relaxation (no $\hat\sigma_z$ channel) at
zero temperature already gives $T_2=2T_1$. Why can $T_2$ never
exceed $2T_1$?
\begin{solution}
From \cref{sec:relaxation}, the coherence with only
$\hat\sigma_-$ decays at $\gamma_\downarrow/2=1/(2T_1)$, so
$T_2=2T_1$. Any additional dephasing channel adds a non-negative
$1/T_\phi$ to $1/T_2$ via \cref{eq:t2-relation}, which can only
shorten $T_2$. Hence $T_2\leq2T_1$ always, with equality iff
$T_\phi\to\infty$.
\end{solution}
\end{exercise}

\begin{exercise}[Deriving the $T_2$ relation]
From the coherence-decay contributions of relaxation
($1/2T_1$) and pure dephasing ($1/T_\phi$), argue that the total
transverse decay rate is their sum, \cref{eq:t2-relation}. Under
what assumption do the two rates simply add?
\begin{solution}
If the relaxation bath and the dephasing bath are statistically
independent and both Markovian, their contributions to the
coherence-decay exponent add linearly:
$\rho_{01}(t)\propto e^{-t/2T_1}e^{-t/T_\phi}=e^{-t/T_2}$ with
$1/T_2=1/2T_1+1/T_\phi$. The assumption is independence and white
(Markovian) statistics; correlated or non-white noise spoils the
simple additivity.
\end{solution}
\end{exercise}

\begin{exercise}[Ramsey fringe with quasi-static noise]
A Ramsey experiment with detuning $\delta$ measures
$\langle\hat\sigma_x(t)\rangle$. Model quasi-static frequency noise
by averaging $\cos[(\delta+\xi)t]$ over a Gaussian distribution of
$\xi$ with variance $\sigma^2$. Show that the fringe envelope
decays as a Gaussian $e^{-\sigma^2t^2/2}$ and identify $T_2^*$.
\begin{solution}
$\langle\cos[(\delta+\xi)t]\rangle_\xi
=\mathrm{Re}\,e^{i\delta t}\langle e^{i\xi t}\rangle
=\cos(\delta t)\,e^{-\sigma^2t^2/2}$ using the Gaussian
characteristic function. The envelope is Gaussian, not exponential,
with $1/T_2^*=\sigma/\sqrt2$ (defining $T_2^*$ by the $1/e$ point of
the envelope, $\sigma^2 (T_2^*)^2/2=1$). This is the inhomogeneous
free-induction decay of \cref{stmt:t2star}, distinct from the
exponential homogeneous decay.
\end{solution}
\end{exercise}

\begin{exercise}[Why $T_1$ samples $S(\omega_q)$]
Using Fermi's golden rule (\cref{thm:fermi}) for the transverse
coupling $\hbar\hat\sigma_x\hat\lambda_\perp$, argue that the
emission rate is proportional to the noise spectral density
evaluated at $+\omega_q$ and the absorption rate to that at
$-\omega_q$, giving \cref{eq:t1-from-noise}.
\begin{solution}
The golden-rule rate for $\ket1\to\ket0$ is
$\propto|\bra0\hat\sigma_x\ket1|^2 S_\perp(\omega_q)$, the bath
supplying a final density of states / spectral weight at the
emitted energy $+\hbar\omega_q$; the $\ket0\to\ket1$ rate is
$\propto S_\perp(-\omega_q)$, requiring the bath to provide
$\hbar\omega_q$. Their sum is $1/T_1$. The transverse operator
$\hat\sigma_x$ is what connects the two levels, so only the bath
spectrum at $\pm\omega_q$ enters.
\end{solution}
\end{exercise}

\begin{exercise}[Why dephasing samples $S(0)$]
Explain why longitudinal $\hat\sigma_z$ coupling produces dephasing
controlled by the \emph{zero-frequency} noise weight, and why for
$1/f$ noise the resulting decay is non-exponential.
\begin{solution}
$\hat\sigma_z$ commutes with $\op H_S$, so it causes no transitions;
it modulates the splitting $\omega_q\to\omega_q+\lambda_\parallel(t)$.
The accumulated random phase is
$\varphi(t)=\int_0^t\lambda_\parallel(t')\dd t'$, whose variance is
$\propto\int S_\parallel(\omega)\,|\,\text{filter}\,|^2\dd\omega$,
peaked at $\omega\to0$ for a free-induction (no-pulse) filter. For
white noise this gives a linear-in-$t$ exponent (exponential decay,
rate $\tfrac12 S_\parallel(0)$); for $1/f$ noise the diverging
low-frequency weight makes the exponent grow faster than linearly
(roughly $t^2\ln$), producing near-Gaussian, non-exponential decay.
\end{solution}
\end{exercise}

\begin{exercise}[Echo refocusing]
Qualitatively explain how a single $\pi$-pulse inserted halfway
through a free-evolution interval (a Hahn echo) removes the effect
of quasi-static frequency noise and recovers $T_2$ from $T_2^*$.
\begin{solution}
Quasi-static noise gives a constant detuning $\xi$ within a shot,
producing phase $\xi\,t/2$ in the first half. The $\pi$-pulse swaps
$\ket0\leftrightarrow\ket1$, flipping the sign of subsequently
accumulated phase, so the second half contributes $-\xi\,t/2$. The
two cancel at the echo time regardless of $\xi$, removing the
quasi-static contribution and leaving only the faster (non-static)
noise --- recovering the homogeneous $T_2$. Noise varying within
the interval is only partially refocused, which is why CPMG (many
$\pi$-pulses) does better.
\end{solution}
\end{exercise}

\begin{exercise}[Purcell relaxation as a $T_1$ channel]
A qubit dispersively coupled to a lossy readout resonator decays
through the resonator (the Purcell effect) at a rate
$\Gamma_{\rm P}=\kappa\,(g/\Delta)^2$, with $\kappa$ the resonator
linewidth, $g$ the coupling, and $\Delta$ the detuning. Explain
how this fits into the $T_1$ budget and how a Purcell filter
suppresses it.
\begin{solution}
Purcell decay is one additive contribution to the relaxation rate,
$1/T_1=\Gamma_{\rm dielectric}+\Gamma_{\rm qp}+\Gamma_{\rm P}+\cdots$;
it is transverse noise delivered by the resonator's photonic
density of states at $\omega_q$. A Purcell filter is a band-stop
element that suppresses the environmental density of states at
$\omega_q$ while passing the readout frequency, reducing
$S_\perp(\omega_q)$ and hence $\Gamma_{\rm P}$ without spoiling
readout. This is the noise-spectral-density viewpoint of
\cref{stmt:noise-rates} applied to engineering.
\end{solution}
\end{exercise}
